% P19 / 225.8: bounded primary-source treatment; Oxford2012-20 remains separate.
\chapter{Motivated Reasoning: Accuracy, Desired Conclusions and Closure}
\label{sec:p19-motivated-reasoning-and-goals}

A reasoner can want an answer to be correct, want a particular answer to be
correct, want the question settled, or prefer a particular way of settling it.
These concerns need not select the same information or stopping point.
The selected Oxford chapter's public abstract distinguishes desired
conclusions, general properties of conclusions, and preferences for judgment
strategies \cite[abstract]{MoldenHiggins2012MotivatedThinkingMRG}.
The account below develops those distinctions through identified original
experiments and a separately identified theory of closure. It does not
identify the original studies with that complete Oxford chapter.

The question is how a goal changes the selection, evaluation or continuation
of reasoning, not whether a response belongs on a list of biases. The
evidence record of \cref{def:p19-evidence-record} remains necessary.
The observed answer, the task's correctness criterion and the process that
produced the answer are three different objects.
\index{motivated reasoning}\index{accuracy motivation}\index{need for closure}

\section{Distinguish the goal from the reasoning operation}
\label{sec:p19-mrg-goals}

\begin{hbdefinition}{def:p19-mrg-goals}{Conclusion, conclusion property and strategy}
A \emph{directional goal} favors a particular conclusion: that one's project
will succeed, that an accusation is false, or that a valued group lacks an
unflattering characteristic. The conclusion's content matters.

An \emph{accuracy goal} favors getting the specified question right, without
fixing which answer must prevail. A \emph{nonspecific closure goal} favors a
definite, stable answer rather than continued ambiguity. Accuracy concerns
the answer's warrant or correctness; closure concerns reaching and retaining
settlement. Neither is defined by the final confidence rating.

A \emph{strategy preference} concerns how judgment proceeds: for example,
searching eagerly for potentially successful candidates or vigilantly
excluding candidates that may be mistaken. Different strategies can pursue
the same desirable outcome. These definitions describe possible explanatory
targets; they do not establish that a particular person has the target.
\end{hbdefinition}

The distinction suggests several operations to examine. A goal may affect
which hypotheses are generated, which records are requested, how demanding
an evidential standard is, or when a search stops. A process can be extensive
yet selective: elaborate criticism of unwelcome evidence is not the same as
equally elaborate examination of all evidence. Conversely, a short answer
can follow an already available sound method. Time spent is not a direct
measure of either accuracy or impartiality.

The monitoring/control interface in
\cref{def:p19-r86-monitoring-control,def:p19-r86-cue-standing}
separates cues to confidence from their validity and use.
The present interface additionally asks what goal makes a cue, an alternative
or an additional check attractive. These are complementary questions, not
two names for one score.

\section{Directional motivation can select a correct falsification}
\label{sec:p19-mrg-directional}

\begin{hbconvention}{def:p19-mrg-report-notation}{Property letters and reported significance}
In the audit case, \(p,q\) name the antecedent and consequent category
properties, and \(\neg p,\neg q\) deny those properties. In the numerical
study reports, \(p\) is instead the reported significance probability of
the specified test. These are explicitly different types of use of the
same letter: a significance probability is neither a category nor its
negation. The test statistic and its reference model determine that
numerical interpretation; it is not an effect size or a probability that
the theory is true.
\end{hbconvention}

\begin{hbexample}{def:p19-mrg-selection-case}{The same audit under two desired conclusions}
Suppose an audit asserts that every member of category \(p\) has property
\(q\), without exception. Four records display \(p,\neg p,q,\neg q\);
each conceals the other property. The task is to inspect exactly the records
that could expose a violation. The \(p\) record may conceal \(\neg q\);
the \(\neg q\) record may conceal \(p\). The other two records cannot
conceal that combination. This is the exceptionless descriptive-rule
benchmark of \cref{sec:p19-r112-selection}.

A manager who wants the audit claim upheld may inspect favorable examples
and leave the potential exception unread. A manager who wants it rejected
may actively seek that exception. The second search can match the task
criterion despite being directional. An accuracy-oriented audit uses the
same criterion whether the rule flatters or threatens the manager.
The example is a worked contrast of possible searches, not an observation
of managers or evidence that every skeptical person reasons correctly.
\end{hbexample}

The original Wason experiments by Dawson, Gilovich and Regan make this
contrast measurable. Study 1 randomly assigned 44 Cornell undergraduates,
participating for course credit, to test a fabricated early-death rule about
high or low emotional lability after their own classification was made
salient. The instruction required two cards. The authors report correct
selection for 46 percent of the personally threatened group versus
10 percent of the nonthreatened group,
\(\chi^2(1,N=44)=6.72,p<.01\).
Neither rule wording alone nor lability classification alone had a detected
effect \cite[pp.~1381--1383, fig.~1]{DawsonGilovichRegan2002WasonMRG}.

Study 2 used 80 Cornell undergraduates, course credit and self-/other-relevant
positive or negative stereotypes, now allowing any number of card selections.
Published correct-selection proportions were 52 percent for negative
self-relevant rules (\(n=31\)), versus 16 percent for the other conditions
combined (\(n=49\)); the reported contrast was
\(\chi^2(1,N=80)=11.26,p<.001\).
The positive-self and positive-other cells were 14 and 20 percent.
The selected report supplies no confidence intervals for these proportions
\cite[pp.~1383--1385, fig.~2]{DawsonGilovichRegan2002WasonMRG}.

The bounded result concerns card choice, not acceptance of the fabricated
claim, clinical risk or stereotype truth. Personally threatening content,
attention, task interpretation and demand characteristics remain possible
routes to the change; the desired conclusion was inferred from the situation,
not independently measured as a mediating variable. Study 2 also used
participant-generated content and sessions with two or three participants.
It is not a randomized change of group membership or a demonstrated transfer
to abstract reasoning. The authors themselves distinguish successful search
in these tasks from acquiring a general logical skill.

\begin{hbwarning}{def:p19-mrg-direction-boundary}{A search motive is not its justification}
Finding a genuine counterexample defeats the stated exceptionless claim,
regardless of why it was sought. Failing to find one in a selective search
does not establish the claim. A person can instead reject a well-supported
claim by applying unattainable standards only to unwelcome evidence.
The actual search, the governing claim and the quality of the inspected
evidence must therefore remain visible. The descriptive finding is not a
recommendation to threaten people or to select desired conclusions.
\end{hbwarning}

\section{Accuracy and closure can diverge before and after commitment}
\label{sec:p19-mrg-closure}

Kruglanski and Webster's theory distinguishes \emph{urgency}, attaining
settlement promptly, from \emph{permanence}, preserving it once attained.
Before a belief has crystallized, an incoming answer can provide wanted
closure; afterward, competing information can threaten it. Their
seizing/freezing distinction thus predicts different openness
at different stages, rather than universally low responsiveness.
The theory separates a wish for one particular answer from a wish for
any definite answer \cite[pp.~263--267]{KruglanskiWebster1996ClosureMRG}.

\begin{hbexample}{def:p19-mrg-closure-case}{Settlement is not a particular favorable verdict}
An investigator has two plausible explanations for a machine failure:
a faulty sensor and an interrupted power supply. An accuracy aim makes a
discriminating inspection relevant. A directional aim favoring the sensor
explanation makes evidence against the power explanation especially welcome.
A closure aim can favor adopting whichever explanation first appears
sufficiently firm, including the unwelcome one.

Before commitment, a colleague's definite explanation can be attractive.
After commitment to the sensor explanation, the same concern for durable
settlement can make a new power-supply record unattractive. If the first
explanation was correct, retaining it can preserve a correct judgment; if it
was wrong, the same pattern can preserve an error. Continued openness can
also admit misleading information. This is an illustration of the theory's
stage distinction, not a measured failure rate or an instruction always to
delay commitment.
\end{hbexample}

A settled answer may be experienced as well warranted because few competing
explanations were generated, not because those alternatives were tested and
excluded. A demand for accuracy must consequently specify an available
method and a relevant check, rather than merely asking for stronger
confidence. Closure is not identical to low intelligence, a pathology, or
the logical closure of a theory under consequence.

\section{Gain-oriented search and protection against errors}
\label{sec:p19-mrg-strategies}

Promotion and prevention concern strategic emphasis, not simply positive
versus negative wording. Promotion emphasizes attaining desired advances;
prevention emphasizes avoiding failures of a desired standard. An opportunity
can be sought through eager candidate generation or protected through
vigilant screening. Which response is preferable additionally depends on
the evidence and the consequences of different mistakes.

\begin{hbdefinition}{def:p19-mrg-four-outcomes}{Recognition outcomes and criterion}
In an old/new recognition task, a \emph{hit} accepts an actually old item;
a \emph{miss} rejects it. A \emph{false alarm} accepts a new item;
a \emph{correct rejection} rejects it. A tendency to accept more items
changes the balance of these outcomes; it is not itself greater ability
to discriminate old from new.

Below \(\beta\) is Crowe and Higgins's reported response-criterion statistic:
one denotes the neutral value, values below one a tendency toward yes,
and values above one toward no. Their reported sensitivity score
\(d'\) is a different measure. The descriptive scores do not supply a
normative cost for a miss or a false alarm.
\end{hbdefinition}

Crowe and Higgins's Study 2 randomly selected 65 paid Columbia
undergraduates from eligible questionnaire respondents and assigned 13 to
each of five frames. Four crossed promotion/prevention with positive/negative
wording of a performance-contingent liked/disliked subsequent activity;
one made that activity noncontingent. Three experimental trials each presented
20 nonsense words for two seconds, a 20-second filler and 40 recognition
items, half old. Testing had no response deadline.
Reported mean \(\beta\) was .92 for promotion and 1.13 for prevention,
\(F(1,60)=6.9,p=.01\); response means were 1.23 versus 1.40 seconds,
\(F(1,60)=6.2,p<.02\).
Sensitivity instead showed a wording-valence effect,
\(F(1,60)=4.6,p<.04\).
No framing effect on measured mood was detected, and the reported focus
effects survived mood adjustment
\cite[pp.~127--130]{CroweHiggins1997StrategyMRG}.

The distinction matters: the result is not that promotion made memory more
accurate or that prevention merely reflected a bad mood. The selected report
does not give confidence intervals for the mean contrasts. Mood adjustment
does not exclude all unmeasured processes, and a particular incentive framing
is not a measurement of a stable personal disposition. The paper prints a
reversed \(d'\) subtraction while describing larger positive scores as
better discrimination; its reported sensitivity comparisons are retained
as reports, not used here to implement that inconsistent formula.

\begin{hbexample}{def:p19-mrg-threshold-case}{Finding candidates versus certifying them}
A reviewer searching a collection for potentially useful records can seek
not to overlook any plausible candidate. A reviewer certifying that every
returned record satisfies a requirement can instead seek not to accept
an unsupported candidate. On a borderline record their actions can differ
without different perceptual ability. Neither role makes an unsupported
claim true. A workflow may need both broad discovery and strict certification,
with their different outputs and error costs stated. This is a worked
strategy contrast, not a finding about professional reviewers.
\end{hbexample}

\section{Motivation, resources and available methods}
\label{sec:p19-mrg-capacity}

A reasoner may lack a needed inference method, be unable to apply it with
the available time or information, or choose not to apply it under the
current goal. These explanations can coexist. An incentive changes neither
the truth conditions of a rule nor automatically the skill needed to use
them. Similarly, effort can be allocated to constructing support for a
preferred answer rather than to a discriminating check.

Viola and colleagues screened 373 volunteers aged 18--36 and selected 60
near the medians of low, medium and high self-report closure-score groups
(\(21,20,19\)). Five were excluded for poor calibration fits. The task
crossed individually calibrated dot-motion ambiguity with stakes:
72 trials, green/no points versus red/30 points gained or lost.
Reaction time (\(\mathrm{RT}\)) was the effort proxy. A promised competitive
ten-euro bonus was followed by ten-euro payment to everyone.
Reported reward-related RT increases occurred in the medium/high groups
(\(p=.001,.005\)), not the low group (\(p=.702\));
ambiguity-related increases occurred in low/medium groups
(\(p=.004,.001\)), not high (\(p=.400\)).
No RT--accuracy correlation was detected in any of the four conditions
\cite[pp.~4--10, figs.~1--4]{ViolaEtAl2015ClosureMRG}.

Closure scores were not randomized. Selection, scale interpretation and
unmeasured participant differences constrain a causal claim about closure.
The red/green stakes cue also changed perceptual context. Longer RT can
reflect a criterion, information accumulation or motor processes, not only
willingness to expend a finite cognitive resource. Nonsignificant
RT--accuracy correlations do not prove the absence of a speed--accuracy
relation. Numerical confidence intervals and a complete post-exclusion
group breakdown are not supplied in the selected report.
Its printed outcome-relevance partial effect-size value, 1.141, exceeds the
allowed range of partial \(\eta^2\); it is not adopted as a valid size.
Repeated printed degrees of freedom for three-group terms also require
clarification before quantitative reanalysis.

\begin{hbmethod}{def:p19-mrg-process-diagnosis}{Separate explanations of the same response}
For a claimed motivational effect, record the question, benchmark, preferred
outcome, supplied information and available strategy. Observe which
alternatives are generated, which information is acquired and when
commitment occurs. Compare the same checks under favorable and unfavorable
content, retaining order, familiarity and instructions.

For a capacity claim, establish whether the method can be used under
adequate time, understandable instructions and available information.
For a resource claim, identify the actual restriction rather than inferring
it from one fast error. For strategic selection, test whether changing
the goal changes information use while those capacities and inputs are held
fixed. An answer score, elapsed time or self-report alone cannot distinguish
these explanations. This is an interpretation procedure, not an assertion
that the three studies isolated every proposed mechanism.
\end{hbmethod}

\section{Consultation without a universal debiasing prescription}
\label{sec:p19-mrg-use}

\begin{hbmethod}{def:p19-mrg-use}{Use the goal distinction on a reasoning record}
Identify whether the aim concerns a specific conclusion, accuracy,
settlement, or a way of proceeding. Distinguish search from evaluation
and evaluation from acceptance. Name the criterion of success: an
exceptionless-rule test, discrimination score, warranted belief, calibrated
confidence, or practical consequence.

Use the actual design to delimit the inference. Random assignment of a
frame does not randomize a person's trait; a successful thematic selection
does not establish general transfer; a slower response does not establish
more accurate reasoning. Preserve unfavorable and inconclusive contrasts
alongside the favorable ones. State an additional evidential and normative
argument before recommending an intervention or a particular error tradeoff.
\end{hbmethod}

The cases separate a counterexample search that happens to fit a directional
aim, a stable answer sought independently of its content, and a criterion
favoring discovery over erroneous acceptance. A failed use treats all three
as believing what one wants, or attributes a result to low capacity
without checking method availability. Another failed use turns an increase
in time or a group difference on a closure scale into advice about whom to
hire or which conclusion to believe.

\begin{proofboundary}{def:p19-mrg-boundary}{Evidence and source-presentation boundary}
This is a bounded primary-study and conceptual comparison, not a common
dataset, a prevalence estimate or a meta-analysis. The experiments use
different tasks, recruited populations and measures; they do not constitute
independent replications of one effect. The closure theory's broader
reviewed experiments and the selected Oxford 2012 chapter's complete
arguments, figures, notes and references form a separate presentation.
Neither an abstract nor these original studies supplies that complete
chapter. No new mathematical theorem, cognitive algorithm or universal
normative prescription is claimed.
\end{proofboundary}
